\documentclass[11pt,a4paper]{scrartcl}

\usepackage[english]{babel}
\usepackage{fontspec}
\usepackage{microtype}
\usepackage[a4paper,left=23mm,right=58mm,top=25mm,bottom=27mm,marginparwidth=42mm,marginparsep=8mm]{geometry}
\usepackage{marginnote}
\usepackage{graphicx}
\usepackage{booktabs}
\usepackage{xcolor}
\usepackage[hidelinks]{hyperref}

\IfFontExistsTF{Libertinus Serif}{%
  \setmainfont{Libertinus Serif}
  \setsansfont{Libertinus Sans}[Scale=0.94]
}{%
  \setmainfont{Latin Modern Roman}
  \setsansfont{Latin Modern Sans}
}

\definecolor{Muted}{HTML}{606060}

\KOMAoptions{parskip=half}
\setkomafont{disposition}{\sffamily\bfseries}

\newcommand{\sidenote}[1]{\marginnote{\footnotesize\color{Muted}#1}}
\newcommand{\newthought}[1]{\par\medskip\noindent{\scshape #1}}

\begin{document}

{\Huge\bfseries Tufte Editorial\par}
\vspace{0.5\baselineskip}
{\large\color{Muted}Main narrative with an active information margin.\par}

\vspace{2\baselineskip}

\section{Information proximity}
\newthought{The defining idea} is not merely a wide margin. The margin carries supporting
information that benefits from remaining close to the primary discussion.
\sidenote{Use sidenotes for context, citations, definitions, compact examples, or secondary commentary.}

The main text remains calm and linear. Supporting material should not interrupt the reader
unless it is part of the principal argument.

\subsection{Margin discipline}
A margin that is mostly empty should not be retained simply to imitate a reference style.
The architecture is justified only when information proximity improves comprehension.

\sidenote{\textbf{Rule:} the margin supports the narrative; it does not become a second competing document.}

\section{Editorial rhythm}
Prefer a narrow, readable measure, quiet section hierarchy, compact footnotes or sidenotes,
and figures placed near the text that interprets them.

\end{document}
